\section*{الفصل \ref{ch:b3:quantum-formalism} --- صياغة الميكانيك الكمومي}

\begin{solution}{exo:b3:quantum-formalism:1}
(أ) $(4 + 1)/5 = 1$ و $(1 + 1)/2 = 1$. (ب) $\braket\varphi\psi =
(2 - \iu)/\sqrt{10}$؛ و $\braket\psi\varphi = (2 + \iu)/\sqrt{10}$:
أي مرافقان. (ج) $|(2 - \iu)|^2/10 = 1/2$. (د) نحلّ
$\braket\psi\chi = 0$: فنجد $\ket\chi = (\ket1 - 2\iu\ket2)/\sqrt5$.
\end{solution}

\begin{solution}{exo:b3:quantum-formalism:2}
الأولى: هرميتية؛ وقيمها $\pm1$ مع $(\ket1 \pm \ket2)/\sqrt2$. والثانية:
هرميتية؛ وقيمها $\pm1$ مع $(\ket1 \pm \iu\ket2)/\sqrt2$. والثالثة: ليست
هرميتية (فمنقولها المرافق يختلف عنها). والرابعة: هرميتية؛
و $\lambda^2 - 5\lambda + 4 = 0$ يعطي $1$ و $4$، بمتجهتين ذاتيتين
$\big({-(1 - \iu)}\ket1 + \ket2\big)/\sqrt3$ و $\big(\ket1 +
(1 + \iu)\ket2\big)/\sqrt3$.
\end{solution}

\begin{solution}{exo:b3:quantum-formalism:3}
(أ) $\braket\beta\alpha = \cos\beta\cos\alpha + \sin\beta\sin\alpha
= \cos(\beta - \alpha)$: فالاحتمال $\cos^2(\beta - \alpha)$. (ب)
الشدّة $\propto$ عدد الفوتونات: أي مالوس. (ج) كل فوتون تجربة
مستقلة: فالتباين $N\mathcal P(1 - \mathcal P)$ --- والضجيج
نفسه يشهد بوجود الفوتونات. (د) $|\cos\beta + \iu\sin\beta|^2/2
= 1/2$ من أجل كل $\beta$: فالضوء الدائري لا يفضّل أيّ محور
مستعرض.
\end{solution}

\begin{solution}{exo:b3:quantum-formalism:4}
(أ) $(\hat x\hat p - \hat p\hat x)f = -\iu\hbar(xf' - (xf)') =
\iu\hbar f$. (ب) $[\hat x^2, \hat p] = 2\iu\hbar\hat x$؛ و $[\hat x,
\hat p^2] = 2\iu\hbar\hat p$. (ج) نضيف ونطرح $\hat B\hat A
\hat C$. (د) بالتراجع مع (ج): $\iu\hbar\,n\hat p^{n-1}$ ---
فالمؤثر $[\hat x, \cdot]$ يعمل عمل $\iu\hbar\,\partial/\partial p$، وهو
الظلّ الكمومي \omterm{def:b3:hamiltonian-mechanics:poisson}{لقوس بواسون} مع $x$.
\end{solution}

\begin{solution}{exo:b3:quantum-formalism:5}
(أ) $|\braket{45^\circ}{V}|^2 = \tfrac12$، ثم انهيار، ثم
$|\braket{H}{45^\circ}|^2 = \tfrac12$: والمجموع $\tfrac14$. (ب) بأخذهما
بالترتيب $60^\circ$ ثم $30^\circ$ (أي ثلاث خطوات مقدارها $30^\circ$):
$(\cos^230^\circ)^3 = 27/64 \approx 0.42$. (ج) تعبر كل خطوة مقدارها
$\pi/2N$ باحتمال $\cos^2(\pi/2N)$: فالنجاة $[\cos^2(\pi/2N)]^N
= 0.25$ و $0.60$ و $0.88$ من أجل $N = 2, 5, 20$. (د) لمّا $N \to \infty$
تؤول النجاة إلى $1$: فقياسات كثيرة لطيفة تسوق الحالة
عبر $90^\circ$ بلا خسارة --- أي القياس مستعملًا مقود قيادة.
\end{solution}

\begin{solution}{exo:b3:quantum-formalism:6}
(أ) $[\hat A, \hat B] = \begin{pmatrix}0 & -2\\ 2 & 0\end{pmatrix}
\neq 0$: أي غير متوافقتين. (ب) قياس $\hat B$ على $\ket1$: يعطي $+1$
يقينًا، ولا حاجة إلى انهيار؛ ثم يعطي $\hat A$ القيمتين $\pm1$
باحتمال $\tfrac12$ لكل منهما. (ج) بعد $\hat A \to +1$ تكون الحالة
$(\ket1 + \ket2)/\sqrt2$؛ فيُنهيها $\hat B$ إلى $\ket1$ أو
$\ket2$ (باحتمال $\tfrac12$ لكل منهما)؛ وفي الحالتين يعطي $\hat A$ الأخير
$\pm1$ باحتمال $\tfrac12$: أي تناقض باحتمال
$\tfrac12$. (د) المرصودات المتبادلة تتقاسم الحالات الذاتية: فالقياس
الأوسط ما كان ليزعج، ولكان التكرار متفقًا
يقينًا.
\end{solution}

\begin{solution}{exo:b3:quantum-formalism:7}
(أ) $\Delta x = \sigma$ و $\Delta p = \hbar/2\sigma$ (فتحويل فورييه
لغاوسية عرضها $\sigma$ عرضه $1/2\sigma$ في
$k$): والجداء $\hbar/2$ بالضبط. (ب) المساواة في كوشي--شوارتز:
$\hat\beta\ket\psi \propto \hat\alpha\ket\psi$ بنسبة تخيلية
صرفة --- ومن أجل $x, p$ لا تحلّ هذه المعادلة التفاضلية إلا
غاوسيات. (ج) $\Delta p \ge \qty{5.3e-25}{kg.m/s}$: ومنه $E \sim
\Delta p^2/2m_{\text{e}} \approx \qty{1}{eV}$ --- فالذرات آلات
بالإلكترون-فولط لأنها علب بالأنغستروم. (د)
$\Delta v \ge \qty{5e-27}{m/s}$: أي لا شيء، أبدًا.
\end{solution}

\begin{solution}{exo:b3:quantum-formalism:8}
(أ) $[\hat H, \hat x] = -\iu\hbar\hat p/m$ و $[\hat H, \hat p] =
\iu\hbar V'(\hat x)$ يعطيان $\dd\langle x\rangle/\dd t = \langle
p\rangle/m$ و $\dd\langle p\rangle/\dd t = -\langle V'\rangle$. (ب)
تغيير المتغير في الجداء الداخلي يُظهر الهرميتية؛
و $\hat\Pi^2 = \mathbb 1$ يفرض القيمتين الذاتيتين $\pm1$. (ج) $V(-x) =
V(x)$ يجعل $\hat H$ أعمى عن الزوجية؛ فوجب أن تكون الحالة الذاتية غير المنحلّة
حالةً ذاتية للمؤثر $\hat\Pi$ أيضًا: زوجية أو فردية. (د)
الثنائي شبه المنحلّ للبئر المزدوجة المتناظرة: حالة زوجية وأخرى
فردية --- أي زوج انقلاب النشادر، المنفصم بالنفق، والمشعّ
عند \qty{24}{GHz}.
\end{solution}

\begin{solution}{exo:b3:quantum-formalism:9}
(أ) تتقاسم الحالتان $E$: فإعلان الطاقة يترك
احتمالين. (ب) يبدّل $\hat S$ الوسمين: وهو هرميتي، ومربعه
المطابقة، ويتبادل مع $\hat H$ المتناظر؛ وداخل المستوي تكون
حالاته الذاتية $(\ket{1,2} \pm \ket{2,1})/\sqrt2$ بقيمتين ذاتيتين
$\pm1$. (ج) $(E, +)$ و $(E, -)$: أي وسمان وحيدان. (د) الانحلال
عنوان ناقص؛ والتناظر الذي يسبّبه يقدّم
الرقم الناقص أيضًا.
\end{solution}

\begin{solution}{exo:b3:quantum-formalism:10}
(أ) $\Delta E\,\Delta A \ge \tfrac12|\langle[\hat H, \hat
A]\rangle| = \tfrac\hbar2|\dd\langle A\rangle/\dd t|$؛ ثم نقسم. (ب)
الزمن وسيط في النظرية لا مؤثر: فالعلاقة
تحدّ من \emph{سرعة} تطور أيّ شيء قابل للقياس، عند
انتشار طاقي معطى. (ج) $\Delta E \sim \hbar/2\tau = \qty{3.3e-8}{eV}$:
أي عرض خط طبيعي مقداره بضعة ميغاهرتز. (د) التغير المرئي في كل
ثانية لا يقتضي غير $\Delta E \gtrsim \qty{5e-35}{J}$ --- ومن أجل
الطاقات العيانية لا \omterm{def:b3:lagrangian-mechanics:coordinates}{قيد} البتة: فللساعات أن تدقّ.
\end{solution}

\begin{solution}{exo:b3:quantum-formalism:11}
(أ) $\|\hat\alpha\psi\|^2 = \braket{\psi}{\hat\alpha^2\psi} =
\Delta A^2$ بهرميتية $\hat\alpha$. (ب) $\braket{\hat\alpha
\psi}{\hat\beta\psi} = \tfrac12\langle\{\hat\alpha,
\hat\beta\}\rangle + \tfrac12\langle[\hat A, \hat B]\rangle$: فالحدّ
الأول حقيقي (الجزء الهرميتي)، والثاني تخيلي صرف.
(ج) $|z|^2 \ge (\operatorname{Im}z)^2$ يعطي المبرهنة؛ وتقتضي المساواة
متجهتين متناسبتين \emph{وانعدام} متوسط \omterm{def:b3:quantum-formalism:commutator}{المبدّلة}
المضادّة. (د) للحالة الذاتية للمؤثر $\hat A$ انحراف $\Delta A = 0$، فيفرض
$0 \ge \hbar/2$: وهو مستحيل على حالات قابلة للتنظيم. أما الأمواج المستوية
فلا ``تحققه'' إلا بكونها مثاليات غير قابلة للتنظيم خارج
\omterm{def:b3:quantum-formalism:state}{فضاء هيلبرت}.
\end{solution}

\begin{solution}{exo:b3:quantum-formalism:12}
(أ) عند $\omega t = \pi/2$، أي $t = T$. (ب) عند كل فحص تكون
الحالة قد دارت بالمقدار $\pi/2N$؛ فتوجد في $\ket1$ باحتمال
$\cos^2(\pi/2N)$ و\emph{تنهار عائدة} إلى $\ket1$؛ والفحوص
مستقلة، ومن ثمّ الجداء. (ج) $0$ و $0.53$ و $0.88$ و $0.98$:
فالقِدر المراقَبة عن كثب لا تغلي أبدًا. (د) مسلَّمة الإسقاط
(م5): إذ تعيد كل مشاهدة التطور إلى خط انطلاقه.
\end{solution}

\begin{solution}{pb:b3:quantum-formalism:1}
\textbf{1.} الدوران يرسل زوجًا متعامدًا نظاميًا إلى زوج متعامد
نظامي: $\braket{\nu_e}{\nu_\mu} = -\cos\theta\sin\theta +
\sin\theta\cos\theta = 0$.
\textbf{2.} المسلَّمة م6: فالطيران الحرّ مولَّد بالمؤثر $\hat H$، وتتطور
حالاته الذاتية استقلاليًا بأطوار --- مهما يكن الأساس الذي
اختاره الإنتاج.
\textbf{3.} $\ket{\psi(0)} = \ket{\nu_\mu} = -\sin\theta\ket{\nu_1}
+ \cos\theta\ket{\nu_2}$.
\textbf{4.} $\ket{\psi(t)} = -\sin\theta\,\eu^{-\iu E_1t/\hbar}
\ket{\nu_1} + \cos\theta\,\eu^{-\iu E_2t/\hbar}\ket{\nu_2}$.
\textbf{5.} تسقط الأطوار الشاملة من كل $|\braket\cdot\cdot|^2$:
فنبقي $\Delta\phi = (E_2 - E_1)t/\hbar$.
\textbf{6.} $E_2 - E_1 = (m_2^2 - m_1^2)c^4/2E = \Delta m^2c^4/2E$.
\textbf{7.} $\braket{\nu_e}{\psi(t)} = \eu^{-\iu E_1t/\hbar}
\sin\theta\cos\theta\,(\eu^{-\iu\Delta\phi} - 1)$.
\textbf{8.} $|\cdots|^2 = \sin^2\theta\cos^2\theta\,|{\eu^{-\iu
\Delta\phi} - 1}|^2 = \sin^22\theta\,\sin^2(\Delta\phi/2)$.
\textbf{9.} لا مزج، أو لا انفصام كتلي: إذًا لا تذبذب. وأيّ
تذبذب مشاهَد يبرهن على أن $\theta \neq 0$ \emph{وأن} $m_1 \neq
m_2$ --- أي إن النيوترينوات تزن.
\textbf{10.} نعوّض $\Delta\phi = \Delta m^2c^4L/2\hbar cE$؛
ويجعل $L_{\text{تذبذب}}$ الوسيطَ مساويًا $\pi$.
\textbf{11.} احتمالا النكهتين مربعا مركبتين في
أساس متعامد نظامي لحالة منظَّمة: فمجموعهما $1$ ---
إذ حفظت واحدية التطور النظيم.
\textbf{12.} لا يُرصد إلا الطور \emph{النسبي}، وهو
يحوي $E_2 - E_1 \propto m_2^2 - m_1^2$: فتتقاصّ الكتل المطلقة.
\textbf{13.} المقدار $\dfrac{L}{4\hbar cE}$ بالوحدات المذكورة:
$\num{e3}\,\unit{m}/(4 \times \num{1.973e-7}\,\unit{eV.m} \times
\num{e9}) = 1.27$ لكل $\unit{eV^2}$.
\textbf{14.} مكتمل النموّ في الأسفل: $1.27\,\Delta m^2 \times 12800
\gtrsim 1$، أي $\Delta m^2c^4 \gtrsim \num{6e-5}\,\unit{eV^2}$؛
ومهمل في الأعلى: $1.27\,\Delta m^2 \times 15 \ll 1$، أي $\ll
\num{5e-2}$: فهو إذًا نحو $10^{-4}$--$10^{-2}\,\unit{eV^2}$.
\textbf{15.} $L_{\text{تذبذب}} = \pi E/(1.27\,\Delta m^2c^4) \approx
\qty{990}{km}$ عند \qty{1}{GeV}: فالمسافة \qty{15}{km} لا تُمسّ، و
\qty{12800}{km} ثلاثة عشر طولًا كاملًا --- أي النقش المشاهَد
بالضبط.
\textbf{16.} على أطوال كثيرة وانتشار في الطاقات يكون
$\langle\sin^2\rangle = \tfrac12$: فالكبت $\tfrac12\sin^22
\theta$؛ والقيمة المقيسة، وهي النصف، تفرض $\sin^22\theta \approx 1$ و
$\theta \approx 45^\circ$ --- فالطبيعة اختارت المزج الأعظمي.
\textbf{17.} $m_2c^2 \approx \sqrt{\num{2.5e-3}} = \qty{0.05}{eV}$:
أي أخفّ عشرة ملايين مرة من الإلكترون --- وهي أخفّ مادة
معروفة، ولا أحد يعرف بعد لماذا.
\textbf{18.} دايا باي: $1.27 \times \num{2.5e-3} \times 1.5/0.004
\approx 1.2$ --- أي من رتبة الواحد، فهي موالَفة؛ وكام-لاند: $1.27 \times
\num{7.5e-5} \times 180/0.004 \approx 4.3$ --- أي من رتبة الواحد من أجل
الانفصام الشمسي: فكل مسافة قاعدية مقياس تداخل مضبوط على
$\Delta m^2$ واحد.
\textbf{19.} المسألة: لم يصل إلا ثلث ما تنبّئ به الشمس من $\nu_e$.
والحلّ: أن الثلثين المفقودين يصلان بنكهات
أخرى، دار إليها $\nu_e$ (وقد أحصت تجربة سنو
المجموع فوجدت الشمس بريئة).
\textbf{20.} يقتضي التذبذب $\Delta m^2 \neq 0$: أي إن نيوترينو
واحدًا على الأقل ثقيل --- وهي أول فيزياء مخبرية وراء
النموذج المعياري الأصلي.
\textbf{21.} ترابط الطور على $10^7$ متر يعني أنه لم يتفاعل
عمليًّا \emph{شيء} في الطريق: فمقاطع التفاعل من الصغر بحيث
تلزم كيلوأطنان من الماء لالتقاط حفنة --- ومن هنا
خمسون ألف طن في سوبر-كاميوكاندي.
\textbf{22.} مؤثرات النكهة قطرية في أساس النكهة،
وهو ليس أساس الطاقة: $[\hat H, \text{النكهة}] \neq 0$ ---
فالنكهة ببساطة ليست مقدارًا محفوظًا في الطيران الحر، بينما
الطاقة وكمية الحركة كذلك.
\textbf{23.} يختلف الانفصامان بعامل ثلاثين: فعند أيّ
$L/E$ يكون أحد التذبذبين نشطًا والآخر إما مجمَّدًا وإما موسَّطًا
تمامًا --- فترى كل تجربة جملةً فعّالة ذات مستويين.
\textbf{24.} أساس الإنتاج $\neq$ أساس الانتشار (أي
الدوران $\theta$)؛ وتقدّم المسلَّمة م6 الطورين؛ وتحوّل \omterm{thm:b3:quantum-formalism:postulates}{قاعدة بورن}
الطور النسبي إلى احتمال.
\textbf{25.} تعطي $\theta \approx 45^\circ$ و $\Delta m^2c^4 =
\qty{2.5e-3}{eV^2}$ لنيوترينو ميوني بالجيغا إلكترون-فولط طول تذبذب
قرب \qty{1000}{km}: أي جبر خطي ذي مستويين، محقَّقًا عبر
الكوكب، وجائزة نوبل لاختفاء نصف تدفق.
\end{solution}
